% Drawing 8/01-1 -- folder 8, batch 01, archive page 1 (PDF page 2).
% Transcription: schéma en carré : $V$ descend de $\varphi^*(C)$, et une flèche « $F = 0$ » revient de $\mathbb{D}^*(G)^{(p)}$ vers $\mathbb{D}^*(G)$
% Kind: square diagram with V and F = 0. Method 2 (schematic redraw from the facsimile). Status: ok.
% Notes: phi*(omega) -> phi*(omega^p) with phi*(C) written above the arrow, from which a double stroke with a heavy head V comes down; under the two ends "=" signs and D*(G), D*(G)^(p); an arrow F = 0 goes back from D*(G)^(p) to D*(G). The upper right label is placed as the transcription reads it.
% Coordinates: 1 unit = 1 mm = 1/100 of the crop width (bbox in page fractions: [0.18, 0.47, 0.42, 0.61]).
% Redrawn by Claude (claude-opus-5-5), issue #73, 2026-11-07; not checked by a human; not compiled (no TeX in the session).
% Rights: reproduces part of the work -- behind RENDER_TRANSCRIPTS, never published by itself.
\documentclass[tikz,border=3pt]{standalone}
\usepackage{amsmath,amssymb}
\usetikzlibrary{patterns,arrows.meta}
\tikzset{pencil/.style={black!45}, struck/.style={black!35}, mark/.style={circle,draw,inner sep=1.3pt}, pt/.style={circle,fill,inner sep=0.9pt}, opt/.style={circle,draw,fill=white,inner sep=1.1pt}}
\begin{document}
\begin{tikzpicture}[x=1mm,y=1mm,line cap=round,line join=round,font=\scriptsize]
\node[inner sep=1pt] at (58,-15) {\tiny $\varphi^*(C)$};
\node[inner sep=1pt] at (35,-36) {\tiny $\varphi^*(\omega)$};
\draw[-{Stealth[length=2mm]}] (46,-34) -- (60,-34);
\node[inner sep=1pt] at (80,-33) {\tiny $\varphi^*(\omega^{p})$};
\draw (52,-22) -- (52.5,-45);
\draw (55,-22) -- (55,-45);
\draw[-{Stealth[length=2mm]},very thick] (53.5,-45) -- (53.5,-51);
\node[inner sep=1pt] at (58,-48) {\tiny $V$};
\node[inner sep=1pt] at (30,-47) {\tiny $=$};
\node[inner sep=1pt] at (80,-43) {\tiny $=$};
\node[inner sep=1pt] at (34,-57) {\tiny $\mathbb{D}^*(G)$};
\node[inner sep=1pt] at (84,-55) {\tiny $\mathbb{D}^*(G)^{(p)}$};
\draw[-{Stealth[length=2mm]}] plot[smooth, tension=0.55] coordinates {(66,-62) (55,-64) (40,-64)};
\node[inner sep=1pt] at (57,-72) {\tiny $F=0$};
\end{tikzpicture}
\end{document}
